% !TEX root = main.tex
\section{Conclusion}
To summarize we first started with a reaction diffusion model for waterborne infection of Cholera. Since our model consists of several PDEs we turn to a numerical approach for a solution. Using second order central finite difference formula for the spatial derivative we derived a numerical scheme for our model. Furthermore, we ensured the no-flux boundary conditions were derived using second order forward/backward difference formulas to maintain a theoretical accuracy of $\mathcal{O}(\Delta x^2)$. We presented our initial conditions and biological parameters based on prior academic research in this related field. As we noted in earlier sections, the proven numerical stability and accuracy ensure the physical validity and reliability of our simulated epidemic dynamics. 

\quad A natural direction for future research using this model would be a more in-depth analysis of the biological impacts of the parameters used. While our current results provide valuable insights, exploring a wider parameter space is required to fully understand the disease dynamics. For example, we primarily simulated different choices of the pathogen's advection rate, $v$, which represents the speed of the river flow. During our simulations, all other parameters remained unchanged, including our initial conditions. Modifying other parameters, such as the pathogen death rate, the diffusion rate, or utilizing a closed SIRB model with a fixed human population, would all be excellent areas of focus for future research. Ultimately, our findings indicate a significant increase in the downstream spread and localized accumulation of the pathogen as the river flow speed increases.

\quad You can find my GitHub repo \href{https://github.com/ethanreyes78/phd-rd-solver}{HERE}.